\section{Die beiden Dateien \file{precoding\_codebooks.cpp}}
\label{sec:dispatcher}

Zwei Dateien gleichen Namens mit klar getrennten Rollen: oben die Fassade und die codebook-unabhängigen
Matrizen, unten die eigentliche 3GPP-Type-I-Mathematik.

\subsection{\file{lib/ran/precoding/precoding\_codebooks.cpp}, die Fassade}

Implementiert den öffentlichen Header \file{include/ocudu/ran/precoding/precoding\_codebooks.h}.
\begin{itemize}
  \item Einfache Matrizen ohne PMI: \code{make\_single\_port()}, \code{make\_one\_layer\_one\_port()},
    \code{make\_one\_layer\_all\_ports()}, \code{make\_identity()}. Sie braucht der FAPI-Adaptor für SSB, PDCCH,
    CSI-RS und das omnidirektionale PDSCH.
  \item \code{get\_single\_panel\_topology()}: Abbildung von $2\times1$, $2\times2$, $4\times1$ auf die
    Antennentopologie.
  \item Zwei Dispatcher: \code{make\_precoding()} und \code{get\_mimo\_matrix\_from\_pmi()}. Die PMI ist ein
    \code{std::variant}; die Visitor-Structs rufen die passende Implementierung.
\end{itemize}

\begin{lstlisting}[language=C++,title={lib/ran/precoding/precoding\_codebooks.cpp}]
/// Precoding weight matrix calculator from a given PMI.
struct precoding_calculator {
  unsigned nof_layers;

  precoding_weight_matrix operator()(const pmi_two_antenna_port& pmi) const
  {
    switch (nof_layers) {
      case 1: return make_one_layer_two_ports(pmi.pmi);
      case 2: return make_two_layer_two_ports(pmi.pmi);
      default: ocudu_assertion_failure("Unsupported PMI codebook configuration");
    }
    return {};
  }

  precoding_weight_matrix operator()(const pmi_typeI_single_panel& pmi) const
  {
    return make_type1_sp_mode1(pmi, nof_layers);
  }

  precoding_weight_matrix operator()(const pmi_typeII& pmi) const { return make_type2(pmi, nof_layers); }
};

precoding_weight_matrix ocudu::make_precoding(const precoding_matrix_indicator& pmi, unsigned nof_layers)
{
  return std::visit(precoding_calculator{nof_layers}, pmi);
}
\end{lstlisting}

Hier wird nichts gerechnet, nur verteilt.

\subsection{\file{lib/ran/precoding/type1\_sp/precoding\_codebooks.cpp}, die Umsetzung}

Die Implementierung von TS\,38.214 Abschnitt 5.2.2.2.1 für ein Single Panel:
\begin{itemize}
  \item \code{make\_layer\_type1\_sp\_mode1()}: ein Steuervektor $[v_{l,m};\,\varphi v_{l,m}]$ für einen Layer.
  \item \code{get\_k1\_and\_k2\_2\_layers()}, \code{get\_k1\_and\_k2\_3\_4\_layers()}: Tabellen 5.2.2.2.1-3 und -4.
  \item \code{make\_type1\_sp\_mode1()}: die vollständige Matrix Ports $\times$ Layer für 1 bis 8 Layer.
  \item \code{get\_beams\_from\_pmi()}, \code{calculate\_mimo\_matrix()}: die zerlegte Variante für Beam-Raster.
\end{itemize}

Der zugehörige Header \file{type1\_sp/precoding\_codebooks.h} liegt bewusst in \file{lib/} statt in
\file{include/}: Diese Funktionen sind intern, von außen soll nur über \code{make\_precoding()} zugegriffen werden.

\subsection{Warum die Aufteilung}

Der Rest des Systems soll nicht wissen, welcher Codebook-Typ gilt. Er ruft \code{make\_precoding(pmi, nof\_layers)}
und bekommt eine Matrix. Jeder Codebook-Typ hat sein eigenes Unterverzeichnis mit eigener Datei desselben Namens.
Zum Lesen geht man von \file{precoding\_codebooks.cpp:151} in den Visitor und von dort in
\file{type1\_sp/precoding\_codebooks.cpp}. Die dritte Datei, \file{type1\_sp/precoding\_codebook\_helpers.cpp},
liefert die Wertebereiche und Bitbreiten der PMI-Felder nach TS\,38.212 Tabelle 6.3.1.1.2-1, die der FAPI-Adaptor
für seine Schleifen und der CSI-Report-Decoder brauchen.
